\documentclass[11pt]{book}
\usepackage{fontspec}
\setmainfont{Latin Modern Roman}
\setsansfont{Latin Modern Sans}
\usepackage{amsmath,amssymb,mathtools}
\usepackage{unicode-math}
\setmathfont{Latin Modern Math}
\usepackage{xurl}
\usepackage[hidelinks]{hyperref}
\newcommand{\sref}[2]{\hyperlink{source:#1:#2}{[S#2]}}
\newcommand{\cataloguescope}[1]{\paragraph{Mathematical question.}#1}
\begin{document}
\section{Diagonal Ramsey numbers}\label{statement}

\subsection{Definitions and mathematical statement}

\paragraph{Shared notation.}$\mathbb N=\{1,2,\ldots\}$, and $\mathbb Z,\mathbb Q,\mathbb R,\mathbb C$ denote the integers, rationals, reals and complex numbers. Any other conventions are specified locally.


For $k\geq2$, $R(k,k)$ is the least integer $N$ such that every red--blue colouring of the edges of the complete graph $K_N$ contains a monochromatic copy of $K_k$.
\paragraph{Statement recorded in the supplied source.}
Does the real limit
\[L=\lim_{k\to\infty}R(k,k)^{1/k}\]
exist? If it exists, determine its value. \sref{statement}{2}
\subsection{Short English statement}
Do the diagonal Ramsey numbers have a well-defined exponential growth rate, and what is that rate?
\subsection{Sources}
\begin{description}
\item[\hypertarget{source:statement:1}{S1}] ULAM.AI and the UnsolvedMath Contributors, \emph{UnsolvedMath: A Curated Collection of Open Mathematics Problems}, catalogue entry. CC BY 4.0. \url{https://huggingface.co/datasets/ulamai/UnsolvedMath}.
\item[\hypertarget{source:statement:2}{S2}] Erdos, Paul, Diagonal Ramsey numbers, Open Problem Garden, node 351. Complete problem, discussion and bibliography (including mathematical image alt text). \url{https://www.openproblemgarden.org/op/diagonal_ramsey_numbers}.
\end{description}
\label{end:statement}
\end{document}
